% !TEX root = ../main.tex
\section{Results}
\label{sec:results}

Results are organized by research question. Within-developer fixed-effects
models and per-category regressions both cluster standard errors at the
developer level. The full PPE sample has $N = 53{,}502$ actors with architecture
fields (50{,}917 agentic-ready); we report both equal-developer-weighted and
actor-weighted fixed-effects estimates.

%% ─────────────────────────────────────────────
\subsection{RQ1: Agentic Readiness and Adoption}
\label{sec:results:rq1}

Agentic-ready tools are adopted far more than non-agentic-ready tools,
identified among the 343 mixed developers who publish both tool types.
\autoref{tab:eligibility_fe} reports the within-developer fixed-effects
estimates. Within the same developer, agentic-ready tools attract
$\hat\beta = 0.540$ more on the log-plus-one scale (SE $= 0.083$,
$p < 10^{-10}$), i.e., \textbf{71.5\% more on the $(1 + \text{monthly\_users})$
scale} than non-agentic-ready tools under equal-developer weighting. The
actor-weighted estimate is smaller but still large and significant
($\hat\beta = 0.334$, 39.6\%, $p < 10^{-6}$). The divergence between the two estimands indicates the gap is
not uniform across developers: larger portfolios---whose many tools dominate
actor weighting---exhibit smaller within-developer gaps. We treat the
equal-developer-weighted estimate as primary, since it answers the
within-developer question without letting a few large portfolios dominate.

\begin{table}[t]
  \centering
  \caption{Agentic readiness and adoption: within-developer fixed effects.}
  \label{tab:eligibility_fe}
  \small
  \begin{tabular}{lccc}
    \toprule
    Estimator & $\hat\beta$ & exp-1 & 95\% CI \\
    \midrule
    Agentic-ready (equal-developer-weighted) & 0.539 & +71.5\% & [0.38, 0.70] \\
    Agentic-ready (actor-weighted) & 0.334 & +39.6\% & [0.21, 0.46] \\
    \bottomrule
    \multicolumn{4}{l}{\footnotesize Identified among 343 mixed developers (16{,}537 PPE actors), developer-clustered} \\
    \multicolumn{4}{l}{\footnotesize SEs; listing-time controls (age, description, category). Equal-developer-} \\
    \multicolumn{4}{l}{\footnotesize weighted gives each developer equal weight. 95\% CI on the $\hat\beta$ (log) scale.} \\
    \multicolumn{4}{l}{\footnotesize PPML with developer FE (count scale): IRR 8.6, $p=0.012$ (equal-developer-weighted);} \\
    \multicolumn{4}{l}{\footnotesize actor-weighted IRR 2.2, n.s.}
  \end{tabular}
\end{table}

The association is not specific to the monthly-user count. On the secondary
outcome $\log(1 + \texttt{runs\_30d})$---the intensity of tool usage---the
within-developer estimate is $+214\%$ (equal-developer-weighted,
$p < 10^{-17}$) and $+108\%$ (actor-weighted, $p < 10^{-5}$), consistent with
and larger than the breadth result.

The association is robust. Excluding whale developers ($\geq$100 PPE actors) leaves
the equal-developer-weighted estimate essentially unchanged ($\hat\beta = 0.558$,
74.6\%, $p < 10^{-9}$). Adding further listing-time tool-level covariates---log
price, log build count, and the number of pricing events---also leaves the
estimate essentially unchanged (71.5\% to 77.5\%, $p < 10^{-10}$), so the
association is not explained by observable tool-level differences in price or
developer effort. Propensity-score weighting on listing-time covariates
yields estimates of 17\%--24\% depending on the estimator
(\autoref{tab:eligibility_psm}), all positive. These estimates provide directional
triangulation; the
treated-to-control ratio is 19.7:1, and the 2{,}585 non-agentic-ready tools are
systematically older than agentic-ready tools (standardized difference in log
age $-0.90$ before weighting, only partly corrected to $+0.23$ after), with balance diagnostics confirming that the within-developer design remains the better-balanced comparison.

Because ``agentic-ready'' is a conjunction of two architecture choices, the
non-agentic-ready group mixes three configurations. To check whether the gap is
driven by one of them---in particular, by standby mode marking a different,
more enterprise tool type---we decompose the architecture into four cells
(\autoref{tab:architecture_2x2}). Within the same developer, the gap is not
confined to standby: relative to ready (limited, non-standby) tools,
full-permission tools are 50\%--56\% lower in adoption and standby-only tools
24\% lower, so the premium is carried by the permission dimension, with standby
a smaller, secondary contributor.

\begin{table}[t]
  \centering
  \caption{Architecture decomposition: adoption across the four
    permission-by-standby cells. Cells report actor count (mean monthly users).}
  \label{tab:architecture_2x2}
  \small
  \begin{tabular}{lcc}
    \toprule
    & \textbf{Standby} & \textbf{Non-standby} \\
    \midrule
    \textbf{Full permissions} & 57 (4.6) & 1,602 (8.6) \\
    \textbf{Limited permissions} & 926 (10.3) & 50,917 (14.3) \\
    \bottomrule
    \multicolumn{3}{l}{\footnotesize Ready = limited, non-standby (bottom-right). Within-developer, relative to} \\
    \multicolumn{3}{l}{\footnotesize ready: full-permission -50.0\% to -54.9\%; standby-only -24.3\%.}
  \end{tabular}
\end{table}

The association also survives restricting to tools with a full exposure window:
equal-developer-weighted estimates are $+73.0\%$ for tools at least 30 days old,
$+89.2\%$ for tools at least 90 days old, and $+61.6\%$ ($p = 0.002$) for tools
at least 180 days old; replacing $\log(\text{age})$
with age-bin fixed effects leaves the estimate at $+69.5\%$. Across the 343
mixed developers, 65.3\% have a positive within-developer ready--non-ready gap
(median $+0.35$ log points), so the result is not driven by a few developers.

To address the concern that agentic-ready and non-agentic-ready tools are simply
different tool types, we match each ready tool to its most similar non-ready
tool within the same developer (\autoref{sec:methods}). Matching improves
covariate balance (log-age standardized difference $-0.90$ to $-0.30$) and
leaves a positive mean premium of roughly 29\% (developer-clustered bootstrap
95\% CI $[0.06, 0.44]$
log points, versus 72\% unmatched): the adoption premium persists after matching,
though it is attenuated, and the median matched gap is zero with
only 47\% of pairs favoring the ready tool, indicating heterogeneity across
developer--tool pairs.

Because the log-plus-one transform compresses the multiplicative gap near zero
(the median tool has one monthly user), we also fit a Poisson pseudo-maximum-
likelihood model with developer fixed effects, which estimates the gap on the
monthly-user count scale directly. The equal-developer-weighted incidence-rate
ratio is 8.6 ($p = 0.012$), confirming a large count-scale gap among typical
portfolios. The actor-weighted IRR is 2.2 and no longer
statistically significant ($p = 0.268$), reflecting the
disproportionate influence of a few large portfolios whose ready and non-ready
tools are closer in absolute counts. The two magnitudes
describe the same gap at different points of a heavily skewed distribution: the
median tool has one user regardless of readiness, so the log-plus-one estimate
($+40\%$ to $+72\%$) moves the median tool from one to roughly two users. The adoption gap is thus large on the
count scale for most developers, not an artifact of the transform.

\begin{table}[t]
  \centering
  \caption{Propensity-score matching for agentic readiness.}
  \label{tab:eligibility_psm}
  \small
  \begin{tabular}{lcc}
    \toprule
    Estimator & $\hat\beta$ & exp-1 \\
    \midrule
    1:1 nearest-neighbor & 0.192 & +21.2\% \\
    Overlap weighting (ATO) & 0.217 & +24.3\% \\
    IPTW (ATT) & 0.156 & +16.9\% \\
    \bottomrule
    \multicolumn{3}{l}{\footnotesize Listing-time propensity score (age, description length,} \\
    \multicolumn{3}{l}{\footnotesize portfolio size, category). Outcome is log-plus-one monthly users.} \\
    \multicolumn{3}{l}{\footnotesize Treated:control = 19.7:1; 1:1 covers a minority of agentic-ready actors.} \\
    \multicolumn{3}{l}{\footnotesize Balance: standardized log-age difference -0.90 before weighting, +0.23 after.}
  \end{tabular}
\end{table}

%% ─────────────────────────────────────────────
\subsection{RQ2: The Margin of Adoption}
\label{sec:results:rq2}

The association appears on both the extensive and conditional intensive
margins (\autoref{tab:eligibility_margins}). On the extensive margin,
agentic-ready tools are \textbf{16.5 percentage points more likely to have any
monthly users} ($p < 10^{-6}$). Conditional on having users, the intensive-margin
association remains large and significant: agentic-ready tools attract
\textbf{60.2\% more} ($p < 10^{-4}$). Entering the top decile is also more
likely for agentic-ready tools ($+10.1$ pp, $p < 10^{-6}$). The association is
thus visible on both margins, consistent with readiness operating through both
discovery and scaling. This two-part decomposition locates where the association appears;
because readiness
itself predicts crossing the zero threshold, the conditional sample of tools
with any users differs in composition between the two groups, so the
intensive-margin estimate is conditional on that selection. Separating agent-mediated from
developer-mediated pathways requires agent-side data (\autoref{fig:mechanism}).

\begin{table}[t]
  \centering
  \caption{Extensive versus intensive margin for agentic readiness.}
  \label{tab:eligibility_margins}
  \small
  \begin{tabular}{lccc}
    \toprule
    Margin & $\hat\beta$ & SE & $p$ \\
    \midrule
    Extensive (any monthly users) & +16.5 pp & 0.032 & <0.001 \\
    Intensive ($\ln(1+\text{MU})$ given MU $>$ 0) & +0.471 & 0.102 & <0.001 \\
    Top decile (MU $\geq$ 10) & +10.1 pp & 0.019 & <0.001 \\
    \bottomrule
    \multicolumn{4}{l}{\footnotesize Within-developer, equal-developer-weighted, developer-clustered SEs.} \\
    \multicolumn{4}{l}{\footnotesize Extensive and top-decile are linear-probability models (coefficient in} \\
    \multicolumn{4}{l}{\footnotesize percentage points); intensive is log-plus-one (coefficient in log points).}
  \end{tabular}
\end{table}

%% ─────────────────────────────────────────────
\subsection{RQ3: Category Heterogeneity and Adoption Inequality}
\label{sec:results:rq3}

\subsubsection{Category heterogeneity.}

The association is positive across most categories in tool-level regressions
(7 of 15 primary-category and 12 of 18 multi-hot are Bonferroni-significant;
\autoref{tab:eligibility_category}). Under developer fixed effects the Ready $\times$ Category interaction is jointly
significant ($F = 6.74$, $p < 10^{-16}$), confirming category-level heterogeneity.
The gradient reorders relative to the tool-level estimates---Agents and AI lose their tool-level
prominence within developer---so the category variation reflects heterogeneity
rather than support for a specific consumer-type mechanism.

\begin{table}[t]
  \centering
  \caption{Per-category agentic-readiness effect (OLS, tool-level controls).}
  \label{tab:eligibility_category}
  \small
  \begin{tabular}{lrcrc}
    \toprule
    Category & Coef. & (SE) & Sig. & $N$ \\
    \midrule
    News             & $+0.6449$ & (0.095) & ***  & 1,231 \\
    Social Media     & $+0.5503$ & (0.116) & ***  & 5,368 \\
    Travel           & $+0.5383$ & (0.150) & ***  & 1,309 \\
    Videos           & $+0.4696$ & (0.165) & **   & 914 \\
    Automation       & $+0.3810$ & (0.058) & ***  & 5,513 \\
    E-commerce       & $+0.3125$ & (0.095) & ***  & 7,364 \\
    Real Estate      & $+0.3047$ & (0.141) & *    & 2,867 \\
    Other            & $+0.2894$ & (0.096) & ***  & 1,209 \\
    AI               & $+0.2882$ & (0.093) & ***  & 3,900 \\
    SEO              & $+0.2361$ & (0.172) & n.s. & 1,548 \\
    Agents           & $+0.2317$ & (0.122) & n.s. & 610 \\
    Jobs             & $+0.2213$ & (0.189) & n.s. & 3,944 \\
    Lead Generation  & $+0.1498$ & (0.141) & n.s. & 8,239 \\
    Developer Tools  & $+0.1192$ & (0.083) & n.s. & 5,747 \\
    MCP Servers      & $-0.1322$ & (0.075) & n.s. & 372 \\
    \bottomrule
    \multicolumn{5}{l}{\footnotesize $^{***}\,p<0.003$ (Bonferroni); $^{**}\,p<0.01$; $^{*}\,p<0.05$; n.s.\,$p\geq 0.05$.} \\
    \multicolumn{5}{l}{\footnotesize Developer-clustered SEs. Controls: log-builds, portfolio size, pricing events.} \\
    \multicolumn{5}{l}{\footnotesize Multi-hot (any-membership) robustness: 12 of 18 categories Bonferroni-significant, all positive.} \\
    \multicolumn{5}{l}{\footnotesize Within-developer (FE) estimates reorder: Agents $-0.10$ and AI $+0.15$ (n.s.);} \\
    \multicolumn{5}{l}{\footnotesize Developer Tools $+0.81$, Social Media $+0.84$, Automation $+0.71$.}
  \end{tabular}
\end{table}

\subsubsection{Within-group adoption inequality.}

We find no robust evidence that adoption is more concentrated among
agentic-ready tools (\autoref{tab:eligibility_concentration}). The Gini
coefficient (the standard scale-invariant measure) is nearly identical across
the two sets (0.935 versus 0.946), with a bootstrap 95\% CI for the difference
of $[-0.04, 0.04]$, narrow around zero; within-category, agentic-ready tools
have a higher Gini in only 7 of 15 categories. There is no reliable evidence of a top-share difference: the face-value gap
(top-1\% 76.8\% versus 74.3\%) has a bootstrap CI of $[-8.7, 25.8]$ pp, consistent with
both no difference and a large one. The total Theil index is higher among
agentic-ready tools (4.13 versus 3.35), but this reflects the agentic-ready
group's larger size and longer tail, and it is no longer distinguishable once
the agentic-ready group is downsampled to the non-agentic-ready size. We
therefore find no detectable concentration premium in the observed readiness
sets.

\begin{table}[t]
  \centering
  \caption{Concentration: agentic-ready versus non-agentic-ready actors.}
  \label{tab:eligibility_concentration}
  \small
  \begin{tabular}{lcc}
    \toprule
    Metric & Agentic-ready & Non-agentic-ready \\
    \midrule
    Gini coefficient & 0.935 & 0.946 \\
    Top 1\% user share & 76.8\% & 74.3\% \\
    Top 5\% user share & 88.7\% & 88.5\% \\
    Top 10\% user share & 92.5\% & 92.4\% \\
    Theil index (total) & 4.13 & 3.35 \\
    Between-category share & 13.1\% & 13.3\% \\
    \bottomrule
    \multicolumn{3}{l}{\footnotesize Within-category Gini higher for agentic-ready in 7 of 15 categories.} \\
    \multicolumn{3}{l}{\footnotesize Bootstrap 95\% CIs for the ready $-$ non-ready difference: Gini $[-0.04, 0.04]$;} \\
    \multicolumn{3}{l}{\footnotesize top-1\% $[-8.7, 25.8]$ pp; top-5\% $[-5.4, 11.7]$ pp; top-10\% $[-3.5, 7.7]$ pp;} \\
    \multicolumn{3}{l}{\footnotesize Theil $[0.07, 2.13]$ (the Theil gap reflects the larger, longer-tailed ready group).} \\
    \multicolumn{3}{l}{\footnotesize Sample-size-equalized (ready downsampled to 2{,}585): Gini CI $[0.85, 0.97]$,} \\
    \multicolumn{3}{l}{\footnotesize top-1\% $[48.3, 89.3]$ pp, Theil $[2.15, 5.33]$---all contain the non-ready values.}
  \end{tabular}
\end{table}
